\chapter{Measurement Regions}%
\label{ch:regions}

The analysis focuses on three kinematic regions: a large phase space
to constrain the interference effects between \ttbar and \tW processes
and improve the modelling over a large kinematic region,
a dedicated phase space at high scales to improve the modelling for future searches,
and a dedicated measurement of the \WbWb production cross section where
the \Wboson boson is directly reconstructed.
The bulk interference region, defined in \cref{sec:regions-bulk}, aims for
the loosest possible selection in the \WbWb phase-space and focuses
on observables sensitive to the interference. The search-like interference
region, defined in \cref{sec:regions-search}, applies additional selection criteria
inspired by BSM searches aiming for as large as possible sensitivity
to the interference. Finally the \Wboson-hadronic measurement region, presented in \cref{sec:regions-wbwb}, focuses on the
overall modelling of this final state and directly considers the \Wboson bosons and \bjets for the definition of the phase space,
which requires the reconstruction of the \Wboson boson (without constraint on intermediate top quarks).
All the region share the same baseline cuts and the summary of the event selection is in \cref{tab:RegionDefinitions}.
The cutflow tables for these selections are available in \cref{app:cutflows}.

\begin{table}[htbp]
  \begin{center}
    \begin{tabular}{l c c c}
      \toprule
      Selection                             & Bulk Interference                      & \Wboson-hadronic     & Search-like Interference         \\
      \midrule
      \( N \)(central jet)                  & \( \geq 4\)                            & \( \geq 3 \)       & \( \geq 4 \)                     \\
      \( N_{\bjets}\)                       & \( \geq 2 \) & \( \geq 2 \) & \( \geq 2 \)                                                               \\
      \( \met + \mtW \)                     & \( (60, \infty) \)  & \( (60, \infty) \)  & \( (60, \infty) \)                                                       \\
      \( \pt(\ell) \)                       & \( (30, \infty) \) & \( (30, \infty) \) & \( (30, \infty) \)                                 \\
      \( \pt(j)\)                           & \( (25, \infty) \) & \( (25, \infty) \) & \( (25, \infty) \)                                       \\
      \( m(\Whad) \)                        & ---                                    & \( (60, 100) \)    & ---                              \\
      \( \pt(\Whad) \)                      & ---                                    & \( (80, \infty) \) & ---                              \\
      \met                                  & ---                                    & ---                & \( \geq -0.625\pt(\ell) + 270 \) \\
      \mblmin                               & ---                                    & ---                & \( (150, \infty) \)              \\
      \( \pt(b_{1}) \)                      & ---                                    & ---                & \( (160, \infty) \)              \\
      \( \pt(b_{2}) \)                      & ---                                    & ---                & \( (50, \infty) \)               \\
      \( \Delta R (b_{1}, b_{2}) \)         & ---                                    & ---                & \( (1.4, \pi) \)                 \\
      \( \pt(j_{1}) \)                      & ---                                    & ---                & \( (50, \infty) \)               \\
      \( \pt(j_{2}) \)                      & ---                                    & ---                & \( (50, \infty) \)               \\
      \bottomrule
    \end{tabular}
  \end{center}
  \caption{A summary of the event selection for bulk-like, \Wboson-hadronic, and search-like regions.}%
  \label{tab:RegionDefinitions}
\end{table}


\section{Baseline Cuts}%
\label{sec:BaselineCuts}
The selection is optimised using simulated MC event samples.
All events are required to contain exactly one signal lepton as defined
in \cref{ch:definitions}. This includes  prompt leptonic \( \tau \) decays, which are considered signal in this analysis.
 Events are selected with a single lepton trigger
and a minimum \pt requirement of \qty{30}{\GeV} is placed on the lepton.
Furthermore the expected number of jets in the
signal is 4 (2 jets from top quark decays and 2 from \Wboson boson decays),
so a preselection condition of at least 4 baseline jets and at least 2 \bjets
in the event is applied. For the \Wboson-hadronic selection this is relaxed to 3 jets
as also a single jet can be reconstructed as the hadronically decaying \Wboson boson.
The neutrino from the second \Wboson boson decay introduces \met.
This allows a cut on \( \met + \mtW > \qty{60}{\GeV} \)
to be applied to remove the QCD background, which in this case enters
the selection as a mis-identified or non-prompt lepton.
The \mtW is defined as the transverse mass computed from
the lepton and \met 4-vectors where the \(z\)-component is set to 0.



\section{Interference sensitivity metrics}

To evaluate if an observable is sensitive to the interference between
\tW and \ttbar two metrics have been defined. 
% First a large difference between DR and DS simulated \tW samples is required:
The first metric measures the difference between DR and DS simulated \tW samples, such that:

\begin{equation}
  \dfrac{N_\mathrm{DS}}{N_\mathrm{DR}}
  \label{eq:metric:DRDS}
\end{equation}
where \( N_\mathrm{DR} \)(\( N_\mathrm{DS} \)) is the number of events passing the selection from the single-top DR\,(DS) samples, respectively.  % chktex 36
Additionally, a separation of up to \qty{20}{\%} between the DR and DS samples is observed in the bulk interference region in most of distributions.
In the search-like interference region, a difference of up to \qty{50}{\%} is observed.
As our modelling of the \tW sample differs, measuring the quantities sensitive to it will help tuning the generators.
This metric only leads to phase spaces where the modelling of \tW is inaccurate. Overall any region where the fraction
of \tW contribution is large, should be measured so a second metric is defined as the fraction of \tW in the signal:

\begin{equation}
  \dfrac{N(\tW)}{N(\tW) + N(\ttbar)}\,,
  \label{eq:metric:ratio}
\end{equation}
In the bulk interference region, a separation of up to \qty{20}{\%} is required between DR and DS samples throughout each distribution.
Additionally, in the search-like interference region, a separation of up to \qty{60}{\%} is necessary.
In addition to the two metrics described in \cref{eq:metric:DRDS,eq:metric:ratio} an additional metric is used for optimising the event selection.
This metric is defined as the difference between the purity of DR and DS,
such that:

\begin{equation}
  FOM = \dfrac{N_\mathrm{DR}}{N_\mathrm{total-simulated(DR)}} - \dfrac{N_\mathrm{DS}}{N_\mathrm{total-simulated(DS)}}
  \label{eq:PurityofDRDS}
\end{equation}

where \(N_\mathrm{total-simulated(DR)}\) and \(N_\mathrm{total-simulated(DS)}\) 
are the total of all simulated events including single-top DR and DS, respectively.  % chktex 36
This metric was used to quantitatively optimise search-like region while keeping sufficient statistics.


\section{Discriminants}
To suppress a large fraction of semi-leptonic \ttbar and SM background,
dominated by \Wjets, a number of discriminating variables have been
selected in addition to the preselection.

\subsubsection{\texorpdfstring{\mblmin}{Minimum invariant mass of a lepton and a b-jet}}
The observable that is most sensitive to the interference
between double and single top production is the invariant mass of
a \bjet and a lepton. An ambiguity is present when pairing the correct
\bjet to the lepton. The solution which provides
the highest sensitivity is:

\begin{equation}
  \mblmin \equiv \min \left\{ m_{b_i \ell} \right\}  % chktex 21
  \label{eq:mblmin}
\end{equation}

where the \(b_i\) goes over all \bjets in the event.
Due to suppression of the doubly resonant contribution at high values
of \mblmin, the differential cross-section above this kinematic endpoint
of \qty{150}{\GeV} has increased sensitivity to interference effects.

\subsubsection{ \( \Delta R (b_{1}, b_{2}) \) }
An additional cut can be introduced on \( \Delta R (b_{1}, b_{2}) \)
to further reduce the \ttbar contribution. It can evade the \mblmin kinematic
bound if the two \btagged jets used in the reconstruction do not originate
from the two top decays, but instead from the \(b\)- and \(c\)-quarks from
a hadronic top decay. In this scenario \( \Delta R (b_{1}, b_{2}) \)
is expected to be small and can be suppressed by requiring
a minimum separation between \btagged jets.

A summary of the quantities used and defined in detail in the following sections is presented in \cref{tab:variables}.

\begin{table}[htbp]
  \begin{center}
    \caption{List of observables used in the three regions of the measurement.}%
    \label{tab:variables}
    \begin{tabular}{|l|l|m{130mm}|}
      \toprule
      \multirow{7}{*}{\rotatebox{90}{\emph{all} }}          & \( \pt(\ell) \)     & Transverse momentum of the lepton                                                                          \\ \cline{2-3}
                                                            & \( \eta(\ell) \)    & Pseudorapidity of the lepton                                                                               \\ \cline{2-3}
                                                            & \( \pt(b_{1,2}) \)  & Transverse momentum of the leading and subleading \bjet                                                    \\ \cline{2-3}
                                                            & \( \pt(j_{1,2}) \)  & Transverse momentum of the leading and subleading jet of any flavour                                       \\ \cline{2-3}
                                                            & \met                & Missing transverse momentum                                                                                \\ \cline{2-3}
                                                            & \mtW                 & Transverse mass between the lepton and the \met                                                            \\ \cline{2-3}
      \midrule
      \multirow{6}{*}{\rotatebox{90}{Bulk and Search-like }} & \amttwo             & Asymmetric stransverse mass~\cite{Lester:1999tx,Barr:2003rg}                                               \\ \cline{2-3}
                                                            & \mblmin             & Minimum invariant mass between the lepton and the \bjet (see text for details)                             \\ \cline{2-3}
                                                            & \mtblmin            & Minimum transverse invariant mass between the \met, the lepton and the \bjet (analogous to \mblmin)        \\ \cline{2-3}
                                                            & \mbb                & Invariant mass of the two \bjets                                                                           \\ \cline{2-3}
                                                            & \ptbb               & Transverse momentum of the 4-vector sum of the two \bjets                                                  \\ \cline{2-3}
                                                            & \deltaRbb           & Spatial separation between 4-vector sum of the two \bjets                                                  \\ \cline{2-3}
      \midrule
      \multirow{12}{*}{\rotatebox{90}{Had W  }}         & \( \pt(\Whad) \)    & Transverse momentum of the hadronic \Wboson candidate (\Whad)                                              \\ \cline{2-3}
                                                            & \( y(\Whad) \)      & Rapidity of the \Whad candidate                                                                            \\ \cline{2-3}
                                                            & \( \pt(j_{OW,1}) \) & Transverse momentum of the leading jet which is not part of the \WbWb system \textit{(outside-of-W-jet)}   \\ \cline{2-3}
                                                            & \( \pt(j_{OW,2}) \) & Transverse momentum of the subleading jet which is not part of the \WbWb system                            \\ \cline{2-3}
                                                            & \mWbbl              & Invariant mass of the \( Wbb\ell \) system                                                                 \\ \cline{2-3}
                                                            & \mbWminimax         & Variation of the min-max variable in Ref.~\cite{TOPQ-2017-05} which uses \Whad and the two leading \bjets:
      \( \min\left( \max(m_{Wb_1},m_{\ell b_2}),\max(m_{\ell b_1},m_{Wb_2})\right)\).
      This definition assigns each of the two \bjets to either \Whad or \( \ell \). This assignment is used for all of the following observables.\                                             \\ \cline{2-3}
                                                            & \ptbl               & Transverse momentum of the 4-vector sum of the lepton and the assigned \bjet                               \\ \cline{2-3}
                                                            & \ptbW               & Transverse momentum of the 4-vector sum of \Whad and the assigned \bjet                                    \\ \cline{2-3}
                                                            & \mbl                & Invariant mass of the 4-vector sum of the lepton and the assigned \bjet                                    \\ \cline{2-3}
                                                            & \mbW                & Invariant mass of the 4-vector sum of \Whad and the assigned \bjet                                         \\ \cline{2-3}
                                                            & \deltaRbl           & Spatial separation between the lepton and the assigned \bjet                                               \\ \cline{2-3}
                                                            & \deltaRbW           & Spatial separation between \Whad and the assigned \bjet                                                    \\ \cline{2-3}
      \bottomrule
    \end{tabular}
  \end{center}
\end{table}



\section{Bulk Interference Region}%
\label{sec:regions-bulk}

On top of the baseline selection only minimal additional cuts have
been investigated. Those are jet multiplicities and baseline jet
momentum selection. Optimisations of the selection are presented in more
detail in \cref{app:optimisation-bulk}. The jet multiplicity selection
stayed at the baseline condition at least 4 jets of which at least 2 should be tagged
as \bjets (pending unfolding optimisations). Several values of baseline
jet \pt have also been tested: \qty{20}{\GeV}, \qty{25}{\GeV}, \qty{30}{\GeV}
and \qty{50}{\GeV}. While minor reduction of uncertainty is expected
for low momenta the tails of distributions sensitive to the interference
do not show a significant difference between individual cut values.
The value of \qty{25}{\GeV} has been chosen as the baseline jet \pt requirement.

Sensitive variables were defined by evaluating the sensitivity metrics
described above on the baseline selection. They can be split into two categories:
basic quantities describing individual particles and more complex compound quantities
where the sensitivity to the \tW and \ttbar interference is targeted.
Out of basic quantities, only lepton \pt and leading \bjet \pt show a high
fraction of \tW events in the signal and a large DR-DS difference,
while the other basic variables do not.
They are shown on \cref{fig:SR_Bulk_Final_l_4j_incl_blinded:basic}.
The yields in the bulk interference region are listed in
\cref{tab:SR_Bulk_Final:yields}.

\begin{table}[htbp]
  \begin{center}
    \input{tables/yields_SR_Bulk_Final_l_4j_incl}
  \end{center}
  \caption{Yields of individual signal and background samples in the
    bulk interference region.
    The total expected yield is using the nominal \tW DR sample.
    Uncertainties are statistical only.}%
  \label{tab:SR_Bulk_Final:yields}
\end{table}

Several compound variables have been defined. Missing transverse momentum is used
as defined in \cref{sec:definitions-met}. The fraction of \tW events is increasing
towards higher values.

The symmetric \(m_{\mathrm{T}2}\) variable is defined as
\begin{equation}
  m_{\mathrm{T}2}(\mathbf{p}_{\mathrm{T},1}, \mathbf{p}_{\mathrm{T},2},\mathbf{E}_{\mathrm{T}}^{\mathrm{miss}}) = \min_{{ \mathbf{q}_{\mathrm{T},1}+\mathbf{q}_{\mathrm{T},2}=\mathbf{E}_{\mathrm{T}}^{\mathrm{miss}}}}[\max(m_{\mathrm{T}}(\mathbf{p}_{\mathrm{T},1},\mathbf{q}_{\mathrm{T},1}),m_{\mathrm{T}}(\mathbf{p}_{\mathrm{T},2},\mathbf{q}_{\mathrm{T},2}))],
\end{equation}
where \(m_{\mathrm{T}}\) is the transverse mass, \(\mathbf{p}_{\mathrm{T},1}\) and  \(\mathbf{p}_{\mathrm{T},2}\) indicate the transverse momentum vectors of two particles of the same mass and \met is the missing transverse momentum vector.
A minimisation is performed over all possible combination of the transverse momentum vectors \(\mathbf{q}_{\mathrm{T},1}\) and \(\mathbf{q}_{\mathrm{T},2}\) that sum up to the missing transverse momentum vector \(\mathbf{E}_{\mathrm{T}}^{\mathrm{miss}}\).
It is used to measure masses of particles pair produced at hadron colliders,
where each particle decays to one particle that is directly observable and another
particle whose existence can only be inferred from missing transverse momenta.
This variable is closely related to the transverse mass variable.
It can be generalised for the asymmetric case where one of the equal-massed particle decays fully visibly
and the other decay involves an invisible particle, as is the case for \ttbar in this analysis~\cite{Lester:2014yga}.
\begin{equation}
  am_{\mathrm{T}2}(\mathbf{p}_{\mathrm{T},b1}, \mathbf{p}_{\mathrm{T},b2},\mathbf{p}_{\mathrm{T},\ell},\mathbf{E}_{\mathrm{T}}^{\mathrm{miss}}) = \min[m_{\mathrm{T}2}(\mathbf{p}_{\mathrm{T},b1}+\mathbf{p}_{\mathrm{T},\ell}, \mathbf{p}_{\mathrm{T},b2},\mathbf{E}_{\mathrm{T}}^{\mathrm{miss}}),m_{\mathrm{T}2}(\mathbf{p}_{\mathrm{T},b1}, \mathbf{p}_{\mathrm{T},b2}+\mathbf{p}_{\mathrm{T},\ell},\mathbf{E}_{\mathrm{T}}^{\mathrm{miss}})]
\end{equation}

A simplified case is just to look at the \mbl variable. As multiple \bjets
are present in the event \mblmin and \mblmax can be defined, where \mblmin
selects the leptonically decaying top quark from the \ttbar process well, while in the \tW case a much
larger fraction is also expected for high mass values. A limitation is a large background contamination.
Similarly \mtblmin and \mtblmax can be defined as combinations of a lepton, \bjet
and \met that give the minimum and maximum values of transverse mass computed
with these objects. While the behaviour of \mtblmin and \mtblmax is similar compared to 
\mblmin and \mblmax, the signal contribution is
smeared more and a better sensitivity due to lower background contamination
is expected for the tail of the distributions.
The \mblmax and \mtblmax variables have lower sensitivity to the interference
and are not used as part of this analysis.


\cref{fig:SR_Bulk_Final_l_4j_incl_blinded:compound} displays the expected distributions
of \met,  \mtW, asymmetric \amttwo, \mtblmin and \mblmin.

Additionally variables related to the \bjet pair are unfolded, namely
\mbb, \ptbb, \deltaRbb and \deltaRbbl. These are shown
in \cref{fig:SR_Bulk_Final_l_4j_incl_blinded:bb}.

More variables, which are not to be unfolded, are shown for their agreement of data and MC are
in \cref{fig:SR_Bulk_Final_l_4j_incl:additional1,fig:SR_Bulk_Final_l_4j_incl:additional2}.

\begin{figure}[htbp]
  \centering
  \subfloat[]{
    \includegraphics[page=28, width=0.37\textwidth]{BulkSelection/SR_Bulk_Final_l_4j_incl_blinded}%
    \label{fig:SR_Bulk_Final_l_4j_incl_blinded:ptl1}
  }
  \subfloat[]{
    \includegraphics[page=1, width=0.37\textwidth]{BulkSelection/SR_Bulk_FinalExtra_l_4j_incl_blinded}%
    \label{fig:SR_Bulk_Final_l_4j_incl_blinded:etal1}
  }
  \\
  \subfloat[]{
    \includegraphics[page=34, width=0.37\textwidth]{BulkSelection/SR_Bulk_Final_l_4j_incl_blinded}%
    \label{fig:SR_Bulk_Final_l_4j_incl_blinded:ptb1}
  }
  \subfloat[]{
    \includegraphics[page=40, width=0.37\textwidth]{BulkSelection/SR_Bulk_Final_l_4j_incl_blinded}%
    \label{fig:SR_Bulk_Final_l_4j_incl_blinded:ptb2}
  }
  \\
  \subfloat[]{
    \includegraphics[page=14, width=0.37\textwidth]{BulkSelection/SR_Bulk_FinalExtra_l_4j_incl_blinded}%
    \label{fig:SR_Bulk_Final_l_4j_incl_blinded:pta1}
  }
  \subfloat[]{
    \includegraphics[page=20, width=0.37\textwidth]{BulkSelection/SR_Bulk_FinalExtra_l_4j_incl_blinded}%
    \label{fig:SR_Bulk_Final_l_4j_incl_blinded:pta2}
  }

  \caption{Basic variables sensitive to the interference between
    \ttbar and single-top.
    \protect\subref{fig:SR_Bulk_Final_l_4j_incl_blinded:ptl1} lepton \pt distribution,
    \protect\subref{fig:SR_Bulk_Final_l_4j_incl_blinded:etal1} lepton \eta distribution,
    \protect\subref{fig:SR_Bulk_Final_l_4j_incl_blinded:ptb1} \( \pt(b_{1}) \) distribution,
    \protect\subref{fig:SR_Bulk_Final_l_4j_incl_blinded:ptb2} \( \pt(b_{2}) \) distribution,
    \protect\subref{fig:SR_Bulk_Final_l_4j_incl_blinded:pta1} \( \pt(j_{1}) \) distribution, and
    \protect\subref{fig:SR_Bulk_Final_l_4j_incl_blinded:pta2} \( \pt(j_{2}) \) distribution.
    In the legend ``Exp'' indicates the prediction with data statistical uncertainty and ``Sys'' includes only MC statistical uncertainty as systematic uncertainty.
    The two predictions differ only in the modelling of the \tW
    contribution and are marked DR (nominal sample) and DS.\@
    The upper ratio panel displays the ratio of the \tW contributions
    for each of the two modelling schemes to the combined \ttbar+\tW
    signal predictions.
    The lower ratio panel displays the ratio of each of the two
    predictions to the nominal one.
    Expected data uncertainty is displayed using error bars.
    Shaded areas show MC statistical uncertainty only.
  }%
  \label{fig:SR_Bulk_Final_l_4j_incl_blinded:basic}
\end{figure}

\begin{figure}[htbp]
  \centering
  \subfloat[]{
    \includegraphics[page=2, width=0.37\textwidth]{BulkSelection/SR_Bulk_Final_l_4j_incl_blinded}%
    \label{fig:SR_Bulk_Final_l_4j_incl_blinded:met}
  }
  \subfloat[]{
    \includegraphics[page=8, width=0.37\textwidth]{BulkSelection/SR_Bulk_FinalExtra_l_4j_incl_blinded}%
    \label{fig:SR_Bulk_Final_l_4j_incl_blinded:mtlep}
  }
  \\
  \subfloat[]{
    \includegraphics[page=8, width=0.37\textwidth]{BulkSelection/SR_Bulk_Final_l_4j_incl_blinded}%
    \label{fig:SR_Bulk_Final_l_4j_incl_blinded:amt2}
  }
  \\
  \subfloat[]{
    \includegraphics[page=22, width=0.37\textwidth]{BulkSelection/SR_Bulk_Final_l_4j_incl_blinded}%
    \label{fig:SR_Bulk_Final_l_4j_incl_blinded:mtblmin}
  }
  \subfloat[]{
    \includegraphics[page=16, width=0.37\textwidth]{BulkSelection/SR_Bulk_Final_l_4j_incl_blinded}%
    \label{fig:SR_Bulk_Final_l_4j_incl_blinded:mblmin}
  }

  \caption{Compound variables sensitive to the interference between
    \ttbar and \tW.\@
    \protect\subref{fig:SR_Bulk_Final_l_4j_incl_blinded:met} \met distribution,
    \protect\subref{fig:SR_Bulk_Final_l_4j_incl_blinded:mtlep}  \mtW distribution,
    \protect\subref{fig:SR_Bulk_Final_l_4j_incl_blinded:amt2} asymmetric \amttwo distribution,
    \protect\subref{fig:SR_Bulk_Final_l_4j_incl_blinded:mtblmin} \mtblmin distribution, and
    \protect\subref{fig:SR_Bulk_Final_l_4j_incl_blinded:mblmin} \mblmin distribution.
    In the legend ``Exp'' indicates the prediction with data statistical uncertainty and ``Sys'' includes only MC statistical uncertainty as systematic uncertainty.
    The two predictions differ only in the modelling of the \tW
    contribution and are marked DR (nominal sample) and DS.\@
    The upper ratio panel displays the ratio of the \tW contributions
    for each of the two modelling schemes to the combined \ttbar+\tW
    signal predictions.
    The lower ratio panel displays the ratio of each of the two
    predictions to the nominal one.
    Expected data uncertainty is displayed using error bars.
    Shaded areas show MC statistical uncertainty only.
  }%
  \label{fig:SR_Bulk_Final_l_4j_incl_blinded:compound}
\end{figure}

\begin{figure}[htbp]
  \centering
  \subfloat[]{
    \includegraphics[page=50, width=0.37\textwidth]{BulkSelection/SR_Bulk_Final_l_4j_incl_blinded}%
    \label{fig:SR_Bulk_Final_l_4j_incl_blinded:mbb}
  }
  \subfloat[]{
    \includegraphics[page=44, width=0.37\textwidth]{BulkSelection/SR_Bulk_Final_l_4j_incl_blinded}%
    \label{fig:SR_Bulk_Final_l_4j_incl_blinded:ptbb}
  }
  \\
  \subfloat[]{
    \includegraphics[page=56, width=0.37\textwidth]{BulkSelection/SR_Bulk_Final_l_4j_incl_blinded}%
    \label{fig:SR_Bulk_Final_l_4j_incl_blinded:drbb}
  }
  \subfloat[]{
    \includegraphics[page=62, width=0.37\textwidth]{BulkSelection/SR_Bulk_Final_l_4j_incl_blinded}%
    \label{fig:SR_Bulk_Final_l_4j_incl_blinded:drbbl}
  }

  \caption{Variables related to two \bjets.\@
    \protect\subref{fig:SR_Bulk_Final_l_4j_incl_blinded:mbb} \mbb distribution,
    \protect\subref{fig:SR_Bulk_Final_l_4j_incl_blinded:ptbb} \ptbb distribution,
    \protect\subref{fig:SR_Bulk_Final_l_4j_incl_blinded:drbb} \deltaRbb distribution, and
    \protect\subref{fig:SR_Bulk_Final_l_4j_incl_blinded:drbbl} \deltaRbbl distribution.
    In the legend ``Exp'' indicates the prediction with data statistical uncertainty and ``Sys'' includes only MC statistical uncertainty as systematic uncertainty.
%
    The two predictions differ only in the modelling of the \tW
    contribution and are marked DR (nominal sample) and DS.\@
    The upper ratio panel displays the ratio of the \tW contributions
    for each of the two modelling schemes to the combined \ttbar+\tW
    signal predictions.
    The lower ratio panel displays the ratio of each of the two
    predictions to the nominal one.
    Expected data uncertainty is displayed using error bars.
    Shaded areas show MC statistical uncertainty only.
  }%
  \label{fig:SR_Bulk_Final_l_4j_incl_blinded:bb}
\end{figure}

\begin{figure}[htbp]
  \centering
  \subfloat[]{
    \includegraphics[page=116, width=0.37\textwidth]{BulkSelection/SR_Bulk_Final_l_4j_incl}%
    \label{fig:SR_Bulk_Final_l_4j_incl:bjet1eta}
  }
  \subfloat[]{
    \includegraphics[page=118, width=0.37\textwidth]{BulkSelection/SR_Bulk_Final_l_4j_incl}%
    \label{fig:SR_Bulk_Final_l_4j_incl:bjet1phi}
  }
  \\
  \subfloat[]{
    \includegraphics[page=120, width=0.37\textwidth]{BulkSelection/SR_Bulk_Final_l_4j_incl}%
    \label{fig:SR_Bulk_Final_l_4j_incl:bjet2eta}
  }
  \subfloat[]{
    \includegraphics[page=122, width=0.37\textwidth]{BulkSelection/SR_Bulk_Final_l_4j_incl}%
    \label{fig:SR_Bulk_Final_l_4j_incl:bjet2phi}
  }
  \\
  \subfloat[]{
    \includegraphics[page=74, width=0.37\textwidth]{BulkSelection/SR_Bulk_Final_l_4j_incl}%
    \label{fig:SR_Bulk_Final_l_4j_incl:lepphi}
  }
  \caption{Variables showing the agreement between data and MC in the bulk interference region:
    \protect\subref{fig:SR_Bulk_Final_l_4j_incl:bjet1eta} leading \bjet \eta distribution,
    \protect\subref{fig:SR_Bulk_Final_l_4j_incl:bjet1phi} leading \bjet \phi distribution,
    \protect\subref{fig:SR_Bulk_Final_l_4j_incl:bjet2eta} subleading \bjet \eta distribution,
    \protect\subref{fig:SR_Bulk_Final_l_4j_incl:bjet2phi} subleading\bjet \phi distribution, and
    \protect\subref{fig:SR_Bulk_Final_l_4j_incl:lepphi} lepton \phi distribution.
    In the legend ``Exp'' indicates the prediction with data statistical uncertainty and ``Sys'' includes only MC statistical uncertainty as systematic uncertainty.
%
    The two predictions differ only in the modelling of the \tW
    contribution and are marked DR (nominal sample) and DS.\@
    The ratio panel displays the ratio of the total MC to the data.
    Expected data uncertainty is displayed using error bars.
    Shaded areas show MC statistical uncertainty only.
  }%
  \label{fig:SR_Bulk_Final_l_4j_incl:additional1}
\end{figure}

\begin{figure}[htbp]
  \centering
  \subfloat[]{
    \includegraphics[page=15, width=0.37\textwidth]{BulkSelection/SR_Bulk_Final_l_4j_incl}%
    \label{fig:SR_Bulk_Final_l_4j_incl:metphi}
  }
  \subfloat[]{
    \includegraphics[page=167, width=0.37\textwidth]{BulkSelection/SR_Bulk_Final_l_4j_incl}%
    \label{fig:SR_Bulk_Final_l_4j_incl:mu}
  }
  \\
  \subfloat[]{
    \includegraphics[page=163, width=0.37\textwidth]{BulkSelection/SR_Bulk_Final_l_4j_incl}%
    \label{fig:SR_Bulk_Final_l_4j_incl:njet}
  }
  \subfloat[]{
    \includegraphics[page=165, width=0.37\textwidth]{BulkSelection/SR_Bulk_Final_l_4j_incl}%
    \label{fig:SR_Bulk_Final_l_4j_incl:nbjet}
  }
  \caption{Variables showing the agreement between data and MC in the bulk interference region:
    \protect\subref{fig:SR_Bulk_Final_l_4j_incl:metphi} \met \phi distribution,
    \protect\subref{fig:SR_Bulk_Final_l_4j_incl:mu} pile-up distribution,
    \protect\subref{fig:SR_Bulk_Final_l_4j_incl:njet} number of jets distribution, and
    \protect\subref{fig:SR_Bulk_Final_l_4j_incl:nbjet} number of \bjets distribution.
    In the legend ``Exp'' indicates the prediction with data statistical uncertainty and ``Sys'' includes only MC statistical uncertainty as systematic uncertainty.
    The two predictions differ only in the modelling of the \tW
    contribution and are marked DR (nominal sample) and DS.\@
    The ratio panel displays the ratio of the total MC to the data.
    Expected data uncertainty is displayed using error bars.
    Shaded areas show MC statistical uncertainty only.
  }%
  \label{fig:SR_Bulk_Final_l_4j_incl:additional2}
\end{figure}

\FloatBarrier


\section{Search-Like Interference Region}%
\label{sec:regions-search}

Besides looking at the bulk selection the aim is to investigate also
a phase space inspired by BSM searches that is particularly sensitive
to the interference between \ttbar and \tW, and where the modelling
approaches can be studied in a region that is particularly
important for BSM searches.
This choice is driven by substantial uncertainties already present
in the search phase space, particularly those pertaining to the theoretical
modelling of single-top events. For example, SUSY searches or searches looking
for Dark Matter~\cite{EXOT-2021-01}. The aim of conducting this study is
to facilitate a more precise calibration of theoretical models related
to single-top processes, ultimately leading to a reduction of associated
uncertainties.
Initial sensitivity analyses and event selection optimizations were carried out,
employing two distinct approaches:
The first utilises a high \met cut, while the second involves a high lepton
\pt cut (see \cref{app:optimisation-searchlike} for details).
The result of these investigations revealed that both variables had a substantial
influence to discriminate difference between the interference
modelling with the DR or DS scheme.\@ Hence, an event selection was
formulated by merging the high \met and high lepton \pt cuts, with the condition
that \met is a linear function of the lepton \pt, such that:

\begin{equation}
  \met > -A\pt\left(\ell\right) + C
\end{equation}

Where \(A\) and \(C\) represent the gradient and the intercept, respectively.
The cut is set at \( \met > -0.625\pt(\ell) + \qty{270}{\GeV} \).
The optimisation of the selection is described in detail
in \cref{app:optimisation-searchlike}.

On top of the baseline selection, other cuts are introduced to discriminate
single-top from \ttbar and \Wjets:
\begin{itemize}
  \item  \( \mblmin \) variable was chosen to discriminate single-top from \( \ttbar \). The cut applied was set at \( \mblmin > \qty{150}{\GeV} \).
  \item \( \Delta R (b_{1}, b_{2}) \) variable was chosen to further reduce the \(\ttbar \) background. The cut applied was set at \( \Delta R (b_{1}, b_{2})  > 1.4\).
  \item Leading \bjet \pt variable was chosen to reduce the \ttbar and \Wjets background. This was set to \( \pt(b_{1}) > \qty{160}{\GeV} \).
  \item 4-jet inclusive selection was preferred over a 3-jet inclusive selection, due to a reduction of the \Wjets background.
  \item 2-\bjet inclusive selection was preferred to a 2-\bjet exclusive selection, due to a reduction of the \Wjets background.
\end{itemize}

The yields in the search-like interference region are listed in
\cref{tab:SR_Searches_Triangular_l_4j_incl:yields}.
Sensitive variables to discriminate predictions using the DR or DS are shown in
\cref{fig:SR_Searches_Triangular_l_4j_incl_blinded:main,fig:SR_Searches_Triangular_l_4j_incl_blinded:compound}.
Furthermore variables related to the two \bjets are shown in
\cref{fig:SR_Searches_Triangular_l_4j_incl_blinded:bb}.

More variables, which are not to be unfolded, are shown for their agreement of data and MC are
in \cref{fig:SR_Searches_Triangular_l_4j_incl:additional1,fig:SR_Searches_Triangular_l_4j_incl:additional2}.

\begin{table}[htbp]
  \begin{center}
    \input{tables/yields_SR_Searches_Triangular_l_4j_incl}
  \end{center}
  \caption{Yields of individual signal and background samples in the
    search-like interference region.
    The total expected yield is using the nominal \tW DR sample.
    Uncertainties are statistical only.}%
  \label{tab:SR_Searches_Triangular_l_4j_incl:yields}
\end{table}

\begin{figure}[htbp]
  \centering
  \subfloat[]{
    \includegraphics[page=21, width=0.37\textwidth]{SearchLikeTriangular/SR_Searches_Triangular_l_4j_incl_blinded}%
    \label{fig:SR_Searches_Triangular_l_4j_incl_blinded:ptl1}
  }
  \subfloat[]{
    \includegraphics[page=1, width=0.37\textwidth]{SearchLikeTriangular/SR_Searches_TriangularExtra_l_4j_incl_blinded}%
    \label{fig:SR_Searches_Triangular_l_4j_incl_blinded:etal1}
  }
  \\
  \subfloat[]{
    \includegraphics[page=27, width=0.37\textwidth]{SearchLikeTriangular/SR_Searches_Triangular_l_4j_incl_blinded}%
    \label{fig:SR_Searches_Triangular_l_4j_incl_blinded:ptb1}
  }
  \subfloat[]{
    \includegraphics[page=33, width=0.37\textwidth]{SearchLikeTriangular/SR_Searches_Triangular_l_4j_incl_blinded}%
    \label{fig:SR_Searches_Triangular_l_4j_incl_blinded:ptb2}
  }
  \\
  \subfloat[]{
    \includegraphics[page=19, width=0.37\textwidth]{SearchLikeTriangular/SR_Searches_TriangularExtra_l_4j_incl_blinded}%
    \label{fig:SR_Searches_Triangular_l_4j_incl_blinded:pta1}
  }
  \subfloat[]{
    \includegraphics[page=25, width=0.37\textwidth]{SearchLikeTriangular/SR_Searches_TriangularExtra_l_4j_incl_blinded}%
    \label{fig:SR_Searches_Triangular_l_4j_incl_blinded:pta2}
  }

  \caption{Variables sensitive to the interference between \ttbar and \tW.\@
    \protect\subref{fig:SR_Searches_Triangular_l_4j_incl_blinded:ptl1} lepton \pt distribution,
    \protect\subref{fig:SR_Searches_Triangular_l_4j_incl_blinded:etal1} lepton \eta distribution,
    \protect\subref{fig:SR_Searches_Triangular_l_4j_incl_blinded:ptb1} \( \pt(b_{1}) \) distribution,
    \protect\subref{fig:SR_Searches_Triangular_l_4j_incl_blinded:ptb2} \( \pt(b_{2}) \) distribution,
    \protect\subref{fig:SR_Searches_Triangular_l_4j_incl_blinded:pta1} \( \pt(j_{1}) \) distribution, and
    \protect\subref{fig:SR_Searches_Triangular_l_4j_incl_blinded:pta2} \( \pt(j_{2}) \) distribution.
    In the legend ``Exp'' indicates the prediction with data statistical uncertainty and ``Sys'' includes only MC statistical uncertainty as systematic uncertainty.
    The two predictions differ only in the modelling of the \tW
    contribution and are marked DR (nominal sample) and DS.\@
    The upper ratio panel displays the ratio of the \tW contributions
    for each of the two modelling schemes to the combined \ttbar+\tW
    signal predictions.
    The lower ratio panel displays the ratio of each of the two
    predictions to the nominal one.
    Expected data uncertainty is displayed using error bars.
    Shaded areas show MC statistical uncertainty only.
  }%
  \label{fig:SR_Searches_Triangular_l_4j_incl_blinded:main}
\end{figure}

\begin{figure}[htbp]
  \centering
  \subfloat[]{
    \includegraphics[page=3, width=0.37\textwidth]{SearchLikeTriangular/SR_Searches_Triangular_l_4j_incl_blinded}%
    \label{fig:SR_Searches_Triangular_l_4j_incl_blinded:met}
  }
  \subfloat[]{
    \includegraphics[page=7, width=0.37\textwidth]{SearchLikeTriangular/SR_Searches_TriangularExtra_l_4j_incl_blinded}%
    \label{fig:SR_Searches_Triangular_l_4j_incl_blinded:mtlep}
  }
  \\
  \subfloat[]{
    \includegraphics[page=13, width=0.37\textwidth]{SearchLikeTriangular/SR_Searches_TriangularExtra_l_4j_incl_blinded}%
    \label{fig:SR_Searches_Triangular_l_4j_incl_blinded:amt2}
  }
  \\
  \subfloat[]{
    \includegraphics[page=9, width=0.37\textwidth]{SearchLikeTriangular/SR_Searches_Triangular_l_4j_incl_blinded}%
    \label{fig:SR_Searches_Triangular_l_4j_incl_blinded:mblmin}
  }
  \subfloat[]{
    \includegraphics[page=15, width=0.37\textwidth]{SearchLikeTriangular/SR_Searches_Triangular_l_4j_incl_blinded}%
    \label{fig:SR_Searches_Triangular_l_4j_incl_blinded:mtblmin}
  }

  \caption{Variables sensitive to the interference between \ttbar and \tW.\@
    \protect\subref{fig:SR_Searches_Triangular_l_4j_incl_blinded:met} \met distribution,
    \protect\subref{fig:SR_Searches_Triangular_l_4j_incl_blinded:mtlep}  \mtW distribution,
    \protect\subref{fig:SR_Bulk_Final_l_4j_incl_blinded:amt2} asymmetric \amttwo distribution,
    \protect\subref{fig:SR_Searches_Triangular_l_4j_incl_blinded:mblmin} \mblmin distribution, and
    \protect\subref{fig:SR_Searches_Triangular_l_4j_incl_blinded:mtblmin} \mtblmin distribution. 
    In the legend ``Exp'' indicates the prediction with data statistical uncertainty and ``Sys'' includes only MC statistical uncertainty as systematic uncertainty.
    The two predictions differ only in the modelling of the \tW
    contribution and are marked DR (nominal sample) and DS.\@
    The upper ratio panel displays the ratio of the \tW contributions
    for each of the two modelling schemes to the combined \ttbar+\tW
    signal predictions.
    The lower ratio panel displays the ratio of each of the two
    predictions to the nominal one.
    Expected data uncertainty is displayed using error bars.
    Shaded areas show MC statistical uncertainty only.
  }%
  \label{fig:SR_Searches_Triangular_l_4j_incl_blinded:compound}
\end{figure}

\begin{figure}[htbp]
  \centering
  \subfloat[]{
    \includegraphics[page=44, width=0.37\textwidth]{SearchLikeTriangular/SR_Searches_Triangular_l_4j_incl_blinded}%
    \label{fig:SR_Searches_Triangular_l_4j_incl_blinded:mbb}
  }
  \subfloat[]{
    \includegraphics[page=38, width=0.37\textwidth]{SearchLikeTriangular/SR_Searches_Triangular_l_4j_incl_blinded}%
    \label{fig:SR_Searches_Triangular_l_4j_incl_blinded:ptbb}
  }
  \\
  \subfloat[]{
    \includegraphics[page=50, width=0.37\textwidth]{SearchLikeTriangular/SR_Searches_Triangular_l_4j_incl_blinded}%
    \label{fig:SR_Searches_Triangular_l_4j_incl_blinded:drbb}
  }
  \subfloat[]{
    \includegraphics[page=56, width=0.37\textwidth]{SearchLikeTriangular/SR_Searches_Triangular_l_4j_incl_blinded}%
    \label{fig:SR_Searches_Triangular_l_4j_incl_blinded:drbbl}
  }

  \caption{Variables related to two \bjets.\@
    \protect\subref{fig:SR_Searches_Triangular_l_4j_incl_blinded:mbb} \mbb distribution,
    \protect\subref{fig:SR_Searches_Triangular_l_4j_incl_blinded:ptbb} \ptbb distribution,
    \protect\subref{fig:SR_Searches_Triangular_l_4j_incl_blinded:drbb} \deltaRbb distribution, and
    \protect\subref{fig:SR_Searches_Triangular_l_4j_incl_blinded:drbbl} \deltaRbbl distribution.
    In the legend ``Exp'' indicates the prediction with data statistical uncertainty and ``Sys'' includes only MC statistical uncertainty as systematic uncertainty.
    The two predictions differ only in the modelling of the \tW
    contribution and are marked DR (nominal sample) and DS.\@
    The upper ratio panel displays the ratio of the \tW contributions
    for each of the two modelling schemes to the combined \ttbar+\tW
    signal predictions.
    The lower ratio panel displays the ratio of each of the two
    predictions to the nominal one.
    Expected data uncertainty is displayed using error bars.
    Shaded areas show MC statistical uncertainty only.
  }%
  \label{fig:SR_Searches_Triangular_l_4j_incl_blinded:bb}
\end{figure}

\begin{figure}[htbp]
  \centering
  \subfloat[]{
    \includegraphics[page=70, width=0.37\textwidth]{SearchLikeTriangular/SR_Searches_Triangular_l_4j_incl}%
    \label{fig:SR_Searches_Triangular_l_4j_incl:bjet1eta}
  }
  \subfloat[]{
    \includegraphics[page=72, width=0.37\textwidth]{SearchLikeTriangular/SR_Searches_Triangular_l_4j_incl}%
    \label{fig:SR_Searches_Triangular_l_4j_incl:bjet1phi}
  }
  \\
  \subfloat[]{
    \includegraphics[page=74, width=0.37\textwidth]{SearchLikeTriangular/SR_Searches_Triangular_l_4j_incl}%
    \label{fig:SR_Searches_Triangular_l_4j_incl:bjet2eta}
  }
  \subfloat[]{
    \includegraphics[page=76, width=0.37\textwidth]{SearchLikeTriangular/SR_Searches_Triangular_l_4j_incl}%
    \label{fig:SR_Searches_Triangular_l_4j_incl:bjet2phi}
  }
  \\
  \subfloat[]{
    \includegraphics[page=46, width=0.37\textwidth]{SearchLikeTriangular/SR_Searches_Triangular_l_4j_incl}%
    \label{fig:SR_Searches_Triangular_l_4j_incl:lepeta}
  }
  \subfloat[]{
    \includegraphics[page=48, width=0.37\textwidth]{SearchLikeTriangular/SR_Searches_Triangular_l_4j_incl}%
    \label{fig:SR_Searches_Triangular_l_4j_incl:lepphi}
  }
  \caption{Variables showing the agreement between data and MC in the search-like region:
    \protect\subref{fig:SR_Searches_Triangular_l_4j_incl:bjet1eta} leading \bjet \eta distribution,
    \protect\subref{fig:SR_Searches_Triangular_l_4j_incl:bjet1phi} leading \bjet \phi distribution,
    \protect\subref{fig:SR_Searches_Triangular_l_4j_incl:bjet2eta} subleading \bjet \eta distribution,
    \protect\subref{fig:SR_Searches_Triangular_l_4j_incl:bjet2phi} subleading \bjet \phi distribution, and
    \protect\subref{fig:SR_Searches_Triangular_l_4j_incl:lepphi} lepton \phi distribution.
    In the legend ``Exp'' indicates the prediction with data statistical uncertainty and ``Sys'' includes only MC statistical uncertainty as systematic uncertainty.
    The two predictions differ only in the modelling of the \tW
    contribution and are marked DR (nominal sample) and DS.\@
    The ratio panel displays the ratio of the total MC to the data.
    Expected data uncertainty is displayed using error bars.
    Shaded areas show MC statistical uncertainty only.
  }%
  \label{fig:SR_Searches_Triangular_l_4j_incl:additional1}
\end{figure}

\begin{figure}[htbp]
  \centering
  \subfloat[]{
    \includegraphics[page=14, width=0.37\textwidth]{SearchLikeTriangular/SR_Searches_Triangular_l_4j_incl}%
    \label{fig:SR_Searches_Triangular_l_4j_incl:metphi}
  }
  \subfloat[]{
    \includegraphics[page=149, width=0.37\textwidth]{SearchLikeTriangular/SR_Searches_Triangular_l_4j_incl}%
    \label{fig:SR_Searches_Triangular_l_4j_incl:mu}
  }
  \\
  \subfloat[]{
    \includegraphics[page=141, width=0.37\textwidth]{SearchLikeTriangular/SR_Searches_Triangular_l_4j_incl}%
    \label{fig:SR_Searches_Triangular_l_4j_incl:njet}
  }
  \subfloat[]{
    \includegraphics[page=145, width=0.37\textwidth]{SearchLikeTriangular/SR_Searches_Triangular_l_4j_incl}%
    \label{fig:SR_Searches_Triangular_l_4j_incl:nbjet}
  }
  \caption{Variables showing the agreement between data and MC in the search-like region:
    \protect\subref{fig:SR_Searches_Triangular_l_4j_incl:metphi} \met \phi distribution,
    \protect\subref{fig:SR_Searches_Triangular_l_4j_incl:mu} pile-up distribution,
    \protect\subref{fig:SR_Searches_Triangular_l_4j_incl:njet} number of jets distribution, and
    \protect\subref{fig:SR_Searches_Triangular_l_4j_incl:nbjet} number of \bjets distribution.
    In the legend ``Exp'' indicates the prediction with data statistical uncertainty and ``Sys'' includes only MC statistical uncertainty as systematic uncertainty.
    The two predictions differ only in the modelling of the \tW
    contribution and are marked DR (nominal sample) and DS.\@
    The ratio panel displays the ratio of the total MC to the data.
    Expected data uncertainty is displayed using error bars.
    Shaded areas show MC statistical uncertainty only.
  }%
  \label{fig:SR_Searches_Triangular_l_4j_incl:additional2}
\end{figure}

\FloatBarrier

\section{\texorpdfstring{\Wboson}{W}-hadronic Measurement Region}%
\label{sec:regions-wbwb}

%Furthermore this analysis focuses also on the \WbWb-process itself.
The purpose of this measurement region is to focus on the \WbWb process
itself and therefore the relevant objects are directly the \Wboson bosons and \bjets, 
instead of lower level objects.
The goal is to measure cross sections as a function of the \WbWb kinematic
observables. These observables include the \pt distribution of the leading \bjet
and of the hadronically decaying \Wboson boson (\Whad). 
In this phase-space, the hadronically decaying \Wboson boson may also be combined with one 
of the \bjets to define properties that are sensitive to doubly-resonant top quark pair contributions.

The phase space is
defined by the baseline cuts outlined in \cref{sec:BaselineCuts}.
In addition, a dedicated reconstruction algorithm for \Whad was developed.
It is explained in detail in \cref{sec:WhadReco}.
In the next step several options of recombining \Whad with one of the
\bjets were investigated. This allowed the definition of a set of observables
similar to those in top-quark pair measurements, and they are expected to be sensitive to the top-quark mass and width. The final phase space definition
 is given at the end of this section, including a list of all the cuts.

\subsection{On the Reconstruction of a W Boson from Light Jets}%
\label{sec:WhadReco}

The set of observables we would like to unfold in the \Wboson-hadronic measurement region (WR)
includes the properties of the hadronically decaying boson \Whad, as well
as combinations of \Whad, \bjets and the lepton.
The explicit reconstruction of the \Whad four-vector and selection based on its properties makes the WR explicitly different to the other two regions.
This requires to derive a reconstruction algorithm for \Whad that can be applied at the detector
(reconstruction) and particle level.
We consider only small-\(R\) jets as the only input to the algorithm, in order to be consistent with the bulk region. 
All light jets passing the baseline \pt and \eta
cuts are considered in the reconstruction of \Whad. 
The study is performed by analysing semi-leptonic \ttbar events.
We require each event to have one hadronically decaying \Wboson and at least one light jet on particle and detector level fulfilling the \(\pt > \qty{25}{GeV}\) and \(|\eta|<2.5\) cuts. No further cuts are applied.

In order to test and compare several recombination schemes,
a criterion on how to measure the matching efficiency has to be defined.
For one event, \(\Delta R\) is calculated between the \Whad candidates
on detector and particle level and the parton level \Wparton boson.
\Wparton denotes the true \Wboson before it decays into quarks.
When \(\Delta R < 0.5\) the candidates are considered as matched.

The \Wboson boson decays into two quarks. The most natural approach is therefore
to reconstruct \Whad from two jets. However, it is also possible for one
of the quarks to radiate a hard gluon. This would result in a third jet.
In contrast, at high \pt(\Wboson), the angle between the two quarks becomes small due to a considerable Lorentz boost.
It is then no longer possible to resolve the two resulting jets, but
the decay products of the \Wboson are clustered into one jet.

\begin{figure}[htpb]
  \centering
  \subfloat[]{
    \includegraphics[width=0.32\textwidth,trim={4cm 0 5cm 0}]{WhadReconstruction/UpdatePlots/WFinder_1j_mWeff}%
    \label{fig:WhadReco_1jm}
  }
  \subfloat[]{
    \includegraphics[width=0.32\textwidth,trim={4cm 0 5cm 0}]{WhadReconstruction/UpdatePlots/WFinder_2j_mWeff}%
    \label{fig:WhadReco_2jm}
  }
  \subfloat[]{
    \includegraphics[width=0.32\textwidth,trim={4cm 0 5cm 0}]{WhadReconstruction/UpdatePlots/WFinder_3j_mWeff}%
    \label{fig:WhadReco_3jm}
  }
  \\
   \subfloat[]{
    \includegraphics[width=0.32\textwidth,trim={4cm 0 5cm 0}]{WhadReconstruction/UpdatePlots/WFinder_1j_ptWeff}%
    \label{fig:WhadReco_1jpt}
  }
  \subfloat[]{
    \includegraphics[width=0.32\textwidth,trim={4cm 0 5cm 0}]{WhadReconstruction/UpdatePlots/WFinder_2j_ptWeff}%
    \label{fig:WhadReco_2jpt}
  }
  \subfloat[]{
    \includegraphics[width=0.32\textwidth,trim={4cm 0 5cm 0}]{WhadReconstruction/UpdatePlots/WFinder_3j_ptWeff}%
    \label{fig:WhadReco_3jpt}
  }
  \caption{Simple reconstruction of a \Wboson boson from light jets.
    Select the 1, 2, or 3-jet combination closest to the \Wboson mass or with
    the highest \pt. This can be done on detector (reconstruction) and
    on particle level. The three curves show the reconstruction efficiency
    between the reconstruction, particle and parton level.
    The efficiency is defined as number of matched events divided by the total
    number of events passing the baseline cuts. An event is considered matched
    between two levels when the respective \Wboson candidates fulfil \( \Delta R < 0.5 \).
    \protect\subref{fig:WhadReco_1jm} Reconstruct from the jet with mass closest to the \Wboson mass,
    \protect\subref{fig:WhadReco_2jm} Reconstruct from the two jet combination with mass closest to the \Wboson mass,
    \protect\subref{fig:WhadReco_3jm} Reconstruct from the three jet combination with mass closest to the \Wboson mass,
    \protect\subref{fig:WhadReco_1jpt} Reconstruct from the jet with highest \pt,
    \protect\subref{fig:WhadReco_2jpt} Reconstruct from the two jet combination with highest \pt,
    \protect\subref{fig:WhadReco_3jpt} Reconstruct from the three jet combination with highest \pt.
  }%
  \label{fig:WhadReco_naive}
\end{figure}

\cref{fig:WhadReco_naive} displays the matching efficiency when the \Wboson
is reconstructed from one, two or three jets. 
%The matching efficiency was evaluated by imposing selections on the light jets themselves. Each light jet is required to have \(\pt > \qty{25}{\GeV}\), be central in the detector (\(|\eta|<2.5\)), and must fail the \(b\)-tagging requirement. This simple selection was used solely to evaluate the matching efficiency for the various algorithms.
In the upper row, the jet
combination whose invariant mass is closest to the \Wboson mass is selected.
The second row selects the hardest jet combination.
It can be observed that \pt ordering is preferred over mass ordering.
The efficiency drops in all panels towards low \pt(\Whad). If the transverse
momentum of the \Wboson boson is small, so is the transverse momentum of its
decay products. Some of the objects originating from the \Wboson boson may not
fulfil the minimal \pt requirement, which makes an accurate reconstruction
impossible. The single jet selection performs especially well at high \pt,
while the two-jet selection breaks down. The reconstruction efficiency of
the three-jet selection is low over the entire \pt range.
It should be noted that the shape of the distribution when matching detector
and particle level (green curve) can vary from the other two curves.
A high detector-particle matching efficiency means that the jets are matched.
It does not mean that the \Wboson boson was reconstructed successfully.

\begin{figure}[htbp]
  \centering
  \subfloat[]{
    \includegraphics[width=0.32\textwidth,trim={4cm 0 5cm 0}]{WhadReconstruction/UpdatePlots/WFinder_1j_ptWeff}%
    \label{fig:WhadReco_1jpt_2}
  }
  \subfloat[]{
    \includegraphics[width=0.32\textwidth,trim={4cm 0 5cm 0}]{WhadReconstruction/UpdatePlots/WFinder_2j_ptWeff}%
    \label{fig:WhadReco_2jpt_2}
  }
  \subfloat[]{
    \includegraphics[width=0.32\textwidth,trim={4cm 0 5cm 0}]{WhadReconstruction/UpdatePlots/WFinder_2j_leadWeff}%
    \label{fig:WhadReco_3jpt_2}
  }
  \\
   \subfloat[]{
    \includegraphics[width=0.32\textwidth,trim={4cm 0 5cm 0}]{WhadReconstruction/UpdatePlots/WFinder_1j_ptWeffCutM}%
    \label{fig:WhadReco_1jpt_m}
  }
  \subfloat[]{
    \includegraphics[width=0.32\textwidth,trim={4cm 0 5cm 0}]{WhadReconstruction/UpdatePlots/WFinder_2j_ptWeffCutM}%
    \label{fig:WhadReco_2jpt_m}
  }
  \subfloat[]{
    \includegraphics[width=0.32\textwidth,trim={4cm 0 5cm 0}]{WhadReconstruction/UpdatePlots/WFinder_2j_leadWeffCutM}%
    \label{fig:WhadReco_3jpt_m}
  }
  \caption{Improved \Wboson selection. The \Wboson boson is reconstructed from the leading jet (left),
    the hardest 2-jet combination (middle) or the combination of leading and subleading jet (right).
    More information on the plots is given in the caption of \cref{fig:WhadReco_naive} and in the text.
    \protect\subref{fig:WhadReco_1jpt_2} Reconstruct \Whad from the jet with highest \pt,
    \protect\subref{fig:WhadReco_2jpt_2} reconstruct from two jet combination with highest \pt,
    \protect\subref{fig:WhadReco_3jpt_2} reconstruct from the combination of leading and subleading jet,
    \protect\subref{fig:WhadReco_1jpt_m} reconstruct from jet with highest \pt including a cut on mass of the reconstructed \(\qty{60}{\GeV} < m(\Whad) < \qty{100}{\GeV}\),
    \protect\subref{fig:WhadReco_2jpt_m} reconstruct from two jet combination with highest \pt including the same mass cut,
    \protect\subref{fig:WhadReco_3jpt_m} reconstruct from leading and subleading jet including the same mass cut.
  }%
  \label{fig:WhadReco_improved}
\end{figure}

Similar plots as in \cref{fig:WhadReco_naive} are shown
in \cref{fig:WhadReco_improved}. The third column shows the efficiency when
the 4-vectors of leading and sub-leading jet are added. The panels in the
lower row include a mass cut on the reconstructed \Wboson candidate
\(\qty{60}{\GeV} < m(\Whad^\mathrm{cand}) < \qty{100}{\GeV} \). 
This increases the matching
efficiency drastically to more than \qty{90}{\%}. The highest efficiency
at medium \pt can be achieved when combining leading and sub-leading jet.
In order to obtain a recombination scheme that is efficient over a large
\pt range it makes sense to pre-select single-jet, two-jet and three-jet
objects. In a next step the \Whad candidate can be chosen from the pre-selected
objects. However, if too many objects are chosen in the pre-selection,
the combinatorial background increases.

After testing several options for
the pre-selection, the matching efficiencies of the best \Whad reconstruction
methods are shown in \cref{fig:WhadReco_best}.
In \cref{fig:WhadReco_5cands} five objects are pre-selected, namely the leading
and the sub-leading jet, the sum of both and the hardest two-jet and three-jet
combinations. The pre-selected object which has a mass closest to the \Wboson
mass is chosen as the \Whad candidate. After the turn-on curve, a plateau
is reached at \(\pt(\Whad) > \qty{200}{\GeV} \). An additional cut on the mass
\(\qty{60}{\GeV} < m(\Whad^\mathrm{cand}) < \qty{100}{\GeV} \) is included in \cref{fig:WhadReco_5cands_mass}. 
This increases the reconstruction efficiency even further to more than \qty{90}{\%} over a large \pt range.
It is even sufficient to pre-select only three candidates, the leading and sub-leading jet as well as the sum of both. 
This can be seen in \cref{fig:WhadReco_5cands_mass}.
The latter selection is applied in the analysis. Therefore, only the leading and sub-leading light jet are considered in the \Whad reconstruction.

The derived reconstruction algorithm can be used to reconstruct \Whad on
detector and on particle level. The reconstruction is easy to understand and
straight-forward to implement. This makes it easier to compare the unfolded
observables in the \Wboson-hadronic region to independent particle level predictions.


\begin{figure}[htbp]
	\centering
	\subfloat[]{
		\includegraphics[width=0.46\textwidth,trim={3cm 0 5cm 0}]{WhadReconstruction/UpdatePlots/WFinder4Weff}%
		\label{fig:WhadReco_5cands}
	}
	\subfloat[]{
		\includegraphics[width=0.46\textwidth,trim={3cm 0 5cm 0}]{WhadReconstruction/UpdatePlots/WFinder4WeffCutM}%
		\label{fig:WhadReco_5cands_mass}
	}
	\\	
	\subfloat[]{
		\includegraphics[width=0.46\textwidth,trim={3cm 0 5cm 0}]{WhadReconstruction/UpdatePlots/WFinder2Weff}%
		\label{fig:WhadReco_3cands}
	}
	\subfloat[]{
		\includegraphics[width=0.46\textwidth,trim={3cm 0 5cm 0}]{WhadReconstruction/UpdatePlots/WFinder2WeffCutM}%
		\label{fig:WhadReco_3cands_mass}
		
	}
	\caption{Efficiencies of the optimized reconstruction of a hadronically
		decaying \Wboson boson from light jets. The highest efficiencies can be
		obtained when a cut on the mass of the reconstructed \Wboson is applied.
		Only events with \(\qty{60} < m(\Whad) < \qty{100}{\GeV}\) are considered in
		\protect\subref{fig:WhadReco_5cands_mass} and \protect\subref{fig:WhadReco_3cands_mass}, while no mass cut is applied
		in \protect\subref{fig:WhadReco_5cands} and \protect\subref{fig:WhadReco_3cands}. 
		Panel~\protect\subref{fig:WhadReco_3cands_mass} shows the efficiency of the method that is used in the analysis. It considers the leading jet, the sub-leading jet, as well as their sum as candidates for the hadronic \Wboson.
		More information is given
		in the caption of \cref{fig:WhadReco_naive} and in the text.}%
	\label{fig:WhadReco_best}
\end{figure}

The \pt resolution for the different \Whad reconstruction algorithms is shown in \cref{fig:WhadResolution}. The resolution plots are obtained from the same \(t\bar{t}\)-sample used in the efficiency study above.
The resolution is defined as
\((\pt(W_\mathrm{reco})-\pt(W_\mathrm{reference}))/\pt(W_\mathrm{reco})\) 
for every event, where the reference \(W_\mathrm{reference}\) is either the true (upper row) or particle level (lower row) \Wboson boson. The plots on the right (\ref{fig:W_res_reco_truth_mass},\ref{fig:W_res_reco_part_mass}) include the mass cut \(60<m(W_\mathrm{reco})<\qty{100}{\GeV}\). To make the comparison easier, the mean \( \mu \) and standard deviation \( \sigma \) are given for each distribution. The values are calculated from the binned histograms.
Applying the mass cut makes the distribution more symmetrical, while at the same time lowering the event yield.
The fraction of events where the leading jet has a mass in the window is expected to be small. It can barely be seen in the figures on the right.
The resolution of combining the two leading jets is very similar to the one with 3 preselected candidates. This is no surprise, since the single jets taken into consideration for the latter algorithm only contribute in the tails, while the resolution distribution is dominated by the low-\pt bulk.

\begin{figure}[htpb]
	\centering
	\subfloat[]{
		\includegraphics[width=0.45\textwidth,trim={2cm 0 2cm 0}]{WhadReconstruction/Resolution/reco_truth}%
		\label{fig:W_res_reco_truth}
	}
	\subfloat[]{
		\includegraphics[width=0.45\textwidth,trim={2cm 0 2cm 0}]{WhadReconstruction/Resolution/reco_truth_mass}%
		\label{fig:W_res_reco_truth_mass}
	}
	\\
	\subfloat[]{
		\includegraphics[width=0.45\textwidth,trim={2cm 0 2cm 0}]{WhadReconstruction/Resolution/reco_part}%
	\label{fig:W_res_reco_part}
	}
	\subfloat[]{
		\includegraphics[width=0.45\textwidth,trim={2cm 0 2cm 0}]{WhadReconstruction/Resolution/reco_part_mass}%
		\label{fig:W_res_reco_part_mass}
	}
	\caption{The \pt resolution of the \Wboson boson for comparing the candidate on detector level (\(W_\mathrm{reco}\)) to the true \Wboson boson (\protect\subref{fig:W_res_reco_truth}, \protect\subref{fig:W_res_reco_truth_mass})
	and the particle level \Wboson 
	(\protect\subref{fig:W_res_reco_part},
	\protect\subref{fig:W_res_reco_part_mass}).
	The mass cut \(60<m(W_\mathrm{reco})<\qty{100}{\GeV}\) is applied in both plots on the right. The mean \( \mu \) and standard deviation \( \sigma \) obtained from the binned histograms are shown in the legend.
	}%
	\label{fig:WhadResolution}
\end{figure}


\subsection{Recombining the hadronic W-boson with a b-jet}%
\label{subsec:wbwb:recombining_had_w}

After reconstructing the \Wboson boson from light jets, one can attempt the
recombination of the \Wboson with one of the \bjets. Only leading and
sub-leading \bjets are considered for this study. Therefore, there are two
possible \Wb combinations. The \Wb system corresponds to a top decaying,
in the case of resonant top quark production.
\cref{fig:TopReconstructionMethods} shows the invariant mass and the \pt
distribution of the \Wb system on truth, particle and detector level.
Six different combination methods are compared in the plots.
These are briefly presented in the following.

\begin{figure}[htbp]
  \centering
  \subfloat[]{
    \includegraphics[width=0.32\textwidth]{TopReconstruction/TOP_MASS_TRUTH}%
    \label{fig:StudyTopMassTruth}
  }
  \subfloat[]{
    \includegraphics[width=0.32\textwidth]{TopReconstruction/TOP_MASS_PART}%
    \label{fig:StudyTopMassPart}
  }
  \subfloat[]{
    \includegraphics[width=0.32\textwidth]{TopReconstruction/TOP_MASS}%
    \label{fig:StudyTopMassReco}
  }
  \\
  \subfloat[]{
    \includegraphics[width=0.32\textwidth]{TopReconstruction/TOP_PT_TRUTH}%
    \label{fig:StudyTopPtTruth}
  }
  \subfloat[]{
    \includegraphics[width=0.32\textwidth]{TopReconstruction/TOP_PT_PART}%
    \label{fig:StudyTopPtPart}
  }
  \subfloat[]{
    \includegraphics[width=0.32\textwidth]{TopReconstruction/TOP_PT}%
    \label{fig:StudyTopPtReco}
  }
  \caption{
    Study on the reconstruction of a hadronically decaying top quark.
    Different selection criteria on how to combine the hadronically decaying
    \Wboson with one of the b jets are compared. Details on the different
    selections are given in the text. The recombined \Wb system is
    denoted with \(t\) in the label of the plots.
    The invariant mass distribution of the reconstructed \Wb system is shown
    on \protect\subref{fig:StudyTopMassTruth} truth,
    \protect\subref{fig:StudyTopMassPart} particle and
    \protect\subref{fig:StudyTopMassReco} detector level, as well as the
    \pt distribution on \protect\subref{fig:StudyTopPtTruth} truth,
    \protect\subref{fig:StudyTopPtPart} particle and
    \protect\subref{fig:StudyTopPtReco} detector level.}%
  \label{fig:TopReconstructionMethods}
\end{figure}

The first method selects the \Wb combination with higher \pt, while
the second one selects the pair that has a mass closer to the top mass.
For the third variation, a parameter \(\chi^2\) is defined which can be
calculated for both combinations \((i=1,j=2\) and \(i=2,j=1)\)
\begin{align*}
    \chi^2_i = {\left( \frac{m(Wb_i)-172.5}{20} \right)}^2 + {\left( \frac{m(lb_j)-100}{30}\right)}^2
\end{align*}
where \(b_1\) and \(b_2\) denote the leading and sub-leading \bjet.
The numerical values in the equation are an estimation of the central
value and width of the \(m(Wb)\) and the \(m(lb)\) distribution, respectively.
The \Wb combination with the smaller \(\chi^2\) is selected as the correct \Wb
combination. The second and third method make explicit use of the top mass.
The fourth method calculates the angular distance \(\Delta R\) between \Wboson
boson and \bjet and selects the pair that is closer.
The fifth methods makes use of the \(m_{Wb}^\mathrm{minmax}\) variable
\begin{align*}
    m_{Wb}^\mathrm{minmax} = \min\left( \max(m(Wb_1),m(lb_2)),\max(m(lb_1),m(Wb_2))\right).
\end{align*}
The observable definition always returns a mass from the first or the second
term. When the observable returns a mass from the first term, \Wboson and
\(b_1\) are combined. When a mass from the second term is returned,
the combination of \(W\) and \(b_2\) is selected.
The last method shown in \cref{fig:TopReconstructionMethods} is a variation
of the \(m_{Wb}^\mathrm{minmax}\) method.
A new observable is defined which also returns a \(Wb\) combination that can
be used as the top candidate
\begin{align*}
    m_{Wb}^\mathrm{maxmin} = \max\left( \min(m(Wb_1),m(Wb_2)),\min(m(lb_1),m(lb_2))\right).
\end{align*}
The width of the mass distribution shown in the upper row of
\cref{fig:TopReconstructionMethods} is increasing when going from the truth
level to particle and detector level. This is expected due to the increasing
combinatorics in the \(W\) reconstruction and the imperfect resolution
on detector level. The selection method choosing the pair closest to the
top mass has the best solution, followed by the minimal \(\chi^2\) method.
The shape of the \pt distribution shown in the lower row is consistent over
all levels. Selecting the harder \Wb pair or the combination with smaller
\(\Delta R(b,\Whad)\) leads to a bias towards higher \pt.

The \(m_{Wb}^\mathrm{minmax}\) and the \(m_{Wb}^\mathrm{maxmin}\) selection
behave similarly. The \pt bias is smaller, the mass resolution is reasonable
and they do not make use of an explicit value for the top mass.
Therefore, the \(m_{Wb}^\mathrm{minimax}\) selection will be used for
the combination of the hadronic \Wboson and a \bjet.

This procedure is used in the definition of several unfolded observables. Specifically, this procedure
is applied to the assignment of the two leading \bjets to the lepton or the reconstructed hadronic \Wboson.
This ensures that the same \bjets cannot be assigned to both objects.
In the text,
observables that have the superscript \textit{selected} use this procedure to select the accompanying \bjets.

In the next step, one can define a matching efficiency for the \Wb
reconstruction. Similarly to what was done for the reconstruction of \Whad,
one can again study the angular distance between the \(Wb\) system on detector
and truth level. A detector level \(Wb\) is considered to be matched
to truth level, when \(\Delta R\) between the two objects is
smaller than \(0.4\).


\begin{figure}[htbp]
  \centering
  \subfloat[]{
    \includegraphics[width=0.4\textwidth]{TopReconstruction/W_PT_BKGD_NODATA}%
    \label{fig:TopMatchingWPtNoData}
  }
  \subfloat[]{
    \includegraphics[width=0.4\textwidth]{TopReconstruction/B_PT_BKGD_NODATA}%
    \label{fig:TopMatchingBPtNoData}
  }
  \caption{Distribution of the transverse momentum of \protect\subref{fig:TopMatchingWPtNoData}
    \Whad and \protect\subref{fig:TopMatchingBPtNoData} the \bjet selected
    with the \(m_{Wb}^\mathrm{minimax}\) method as described in the text.
    ``Correct match'' denotes events where
    \(\Delta R(Wb_\mathrm{reco},Wb_\mathrm{truth}) < 0.4\),
    where \Wb is the sum of the \Wboson and \bjet 4-vectors.
    ``No \Wb'' denotes events where either \Whad or the \bquark does not exist
    on truth level.}%
    \label{fig:TopMatchNoData}
\end{figure}

\cref{fig:TopMatchNoData} shows the detector level \Wboson and \bjet \pt distributions. 
The \WbWb phase space selection is applied on detector level.
It includes a tighter cut on the \Wboson mass
\((60 < m(\Whad) < \qty{100}{\GeV} \mathrm{and~} \pt > \qty{80}{\GeV})\).
These stricter cuts help to improve the matching efficiency. The plots show the fraction of correctly matched events for the \mbWminimax method.
%All six combination schemes were tested.
%The best results were achieved with the \(m_{Wb}^\mathrm{minimax}\) method.
%Therefore, the results of the other methods are not shown here.
Only the \bjets that were selected using the
\mbWminimax method are considered.
The background contribution is small and can not be seen on a linear scale.
The events of the signal processes are split into three categories.
Correct match and wrong match describe the events where
\(\Delta R({(Wb)}_\mathrm{truth},{(Wb)}_\mathrm{reco}) < 0.4\) is fulfilled or not.
The third category ``No Wb'' denotes the events where either \Whad or \bjet
does not exist on truth level. These are mainly \ttbar dilepton events that
were not discarded in the detector level event selection.
The fraction of ``No Wb'' events is in the low percent range.
The fraction of correctly matched events does not depend on \bjet \pt,
while a lower efficiency at small \Wboson \pt can be observed.
This discrepancy can be traced back to a lower reconstruction efficiency
of the \Whad at lower \pt. This can be seen in \cref{fig:TopMatchingWPtNoData}.
The \Wb matching efficiency could be further improved when increasing
the \Wboson \pt cut.

\begin{figure}[htbp]
  \centering
  \subfloat[]{
    \includegraphics[width=0.4\textwidth]{TopReconstruction/W_MASS_BKGD}%
    \label{fig:TopMatchingWMass}
  }
  \subfloat[]{
    \includegraphics[width=0.4\textwidth]{TopReconstruction/W_PT_BKGD}%
    \label{fig:TopMatchingWPt}
  }
  \\
  \subfloat[]{
    \includegraphics[width=0.4\textwidth]{TopReconstruction/TOP_MASS_BKGD}%
    \label{fig:TopMatchingTopMass}
  }
  \subfloat[]{
    \includegraphics[width=0.4\textwidth]{TopReconstruction/TOP_PT_BKGD}%
    \label{fig:TopMatchingTopPt}
  }
  \caption{
    Mass and \pt distribution of the reconstructed \Whad and \(Wb\) object. The Figure compares the MC prediction to the data. The \(Wb\) combination is found with the \(m_{Wb}^\mathrm{minimax}\) method, as described in the text. More information is given in the caption of~\cref{fig:TopMatchNoData} and in the text.
    }%
    \label{fig:TopMatchWithData}
\end{figure}

The invariant mass distribution of \Whad and \Wb is displayed
in \cref{fig:TopMatchWithData}, as well as the \pt distributions.
The applied detector level cuts on \(m(\Whad)\) and \(\pt(\Whad)\) can be seen
in the upper row. The \(m(Wb)\) distribution peaks at the top mass.
The tails are slightly asymmetrical towards higher masses.
An offset between correctly and wrongly matched events can be seen
in the \(\pt(Wb)\) distribution. The fraction of correctly matched events
can be improved when increasing the \(\pt(\Whad)\) cut,
or by introducing an additional \(\pt(Wb)\) cut.
In the boosted regime \((\pt > \qty{200}{\GeV})\) the matching efficiency
is almost perfect. The agreement between data an Monte Carlo is reasonable.

\subsection{W-hadronic Region Final Event Selection}%

The event selection for the \Wboson-hadronic region follows the strategy outlined
in \cref{sec:BaselineCuts}.
Additional cuts are derived from the \Whad reconstruction studies.
The requirement of at least 4 baseline jets is relaxed.
After selecting exactly 2 \bjets, the remaining jets are passed
on to the \Whad reconstruction algorithm.
\Whad is reconstructed from the one or two leading jets, as described
in \cref{sec:WhadReco}.
The reconstructed \Whad candidate has to fulfil
\begin{itemize}
    \item \( 60 < m(\Whad) < \qty{100}{\GeV} \)
    \item \( \qty{80}{\GeV} < \pt(\Whad) \)
\end{itemize}
Several thresholds for the mass and \pt cut were tested in dedicated studies.
It was found that the stricter cuts chosen in the end improve the \Whad reconstruction.
At the same time they increase the signal to background ration by almost a factor of 2.
The exact values of the \Whad cuts do not affect the size of the systematic uncertainties.

The yields in the \Wboson-hadronic region are listed in
\cref{tab:SR_Whad_Final:yields}.

\begin{table}[htbp]
  \begin{center}
    \input{tables/yields_SR_Whad_Final_l_Whad}
  \end{center}
  \caption{Yields of individual signal and background samples in the
    \Wboson-hadronic measurement region.
    The total expected yield is using the nominal \tW DR sample.
    Uncertainties are statistical only.}%
  \label{tab:SR_Whad_Final:yields}
\end{table}

More variables, which are not to be unfolded, are shown for their agreement of data and MC are
in \cref{fig:SR_Whad_Final_l_Whad_particle:additional1,fig:SR_Whad_Final_l_Whad_particle:additional2}.

\begin{figure}[htbp]
  \centering
  \subfloat[]{
    \includegraphics[page=5, width=0.3\textwidth]{WhadSelection/SR_Whad_FinalExtra_l_Whad_particle}%
    \label{fig:SR_Whad_Final_l_Whad_particle:met}
  }
  \subfloat[]{
    \includegraphics[page=5, width=0.3\textwidth]{WhadSelection/SR_Whad_Final_l_Whad_particle}%
    \label{fig:SR_Whad_Final_l_Whad_particle:metphi}
  }
  \\
  \subfloat[]{
    \includegraphics[page=9, width=0.3\textwidth]{WhadSelection/SR_Whad_Final_l_Whad_particle}%
    \label{fig:SR_Whad_Final_l_Whad_particle:lep1pt}
  }
  \subfloat[]{
    \includegraphics[page=13, width=0.3\textwidth]{WhadSelection/SR_Whad_Final_l_Whad_particle}%
    \label{fig:SR_Whad_Final_l_Whad_particle:lep1eta}
  }
  \subfloat[]{
    \includegraphics[page=15, width=0.3\textwidth]{WhadSelection/SR_Whad_Final_l_Whad_particle}%
    \label{fig:SR_Whad_Final_l_Whad_particle:lep1phi}
  }
  \caption{Variables showing the agreement between data and MC in the \Wboson-hadronic region:
    \protect\subref{fig:SR_Whad_Final_l_Whad_particle:met} \met distribution,
    \protect\subref{fig:SR_Whad_Final_l_Whad_particle:metphi} \met \phi distribution,
    \protect\subref{fig:SR_Whad_Final_l_Whad_particle:lep1pt} lepton \pt distribution,
    \protect\subref{fig:SR_Whad_Final_l_Whad_particle:lep1eta} lepton \eta distribution, and
    \protect\subref{fig:SR_Whad_Final_l_Whad_particle:lep1phi} lepton \phi distribution.
    The four predictions differ only in the modelling of the \tW
    contributions, where the predictions using dynamic scales are marked
    DR (nominal sample) and DS while the legacy samples with fixed scales
    are marked DRF and DSF.\@
    The ratio panel displays the ratio of the total MC to the data.
    Expected data uncertainty is displayed using error bars.
    Shaded areas show MC statistical uncertainty only.
  }%
  \label{fig:SR_Whad_Final_l_Whad_particle:additional1}
\end{figure}

\begin{figure}[htbp]
  \centering
  \subfloat[]{
    \includegraphics[page=25, width=0.23\textwidth]{WhadSelection/SR_Whad_Final_l_Whad_particle}%
    \label{fig:SR_Whad_Final_l_Whad_particle:bjet1pt}
  }
  \subfloat[]{
    \includegraphics[page=36, width=0.23\textwidth]{WhadSelection/SR_Whad_Final_l_Whad_particle}%
    \label{fig:SR_Whad_Final_l_Whad_particle:bjet1eta}
  }
  \subfloat[]{
    \includegraphics[page=38, width=0.23\textwidth]{WhadSelection/SR_Whad_Final_l_Whad_particle}%
    \label{fig:SR_Whad_Final_l_Whad_particle:bjet1phi}
  }
  \\
  \subfloat[]{
    \includegraphics[page=31, width=0.23\textwidth]{WhadSelection/SR_Whad_Final_l_Whad_particle}%
    \label{fig:SR_Whad_Final_l_Whad_particle:bjet2pt}
  }
  \subfloat[]{
    \includegraphics[page=40, width=0.23\textwidth]{WhadSelection/SR_Whad_Final_l_Whad_particle}%
    \label{fig:SR_Whad_Final_l_Whad_particle:bjet2eta}
  }
  \subfloat[]{
    \includegraphics[page=42, width=0.23\textwidth]{WhadSelection/SR_Whad_Final_l_Whad_particle}%
    \label{fig:SR_Whad_Final_l_Whad_particle:bjet2phi}
  }
  \\
  \subfloat[]{
    \includegraphics[page=175, width=0.23\textwidth]{WhadSelection/SR_Whad_Final_l_Whad_particle}%
    \label{fig:SR_Whad_Final_l_Whad_particle:mu}
  }
  \subfloat[]{
    \includegraphics[page=171, width=0.23\textwidth]{WhadSelection/SR_Whad_Final_l_Whad_particle}%
    \label{fig:SR_Whad_Final_l_Whad_particle:njets}
  }
  \subfloat[]{
    \includegraphics[page=173, width=0.23\textwidth]{WhadSelection/SR_Whad_Final_l_Whad_particle}%
    \label{fig:SR_Whad_Final_l_Whad_particle:nbjets}
  }
  \caption{Variables showing the agreement between data and MC in the \WbWb region:
    \protect\subref{fig:SR_Whad_Final_l_Whad_particle:bjet1pt} leading \bjet \pt distribution,
    \protect\subref{fig:SR_Whad_Final_l_Whad_particle:bjet1eta} leading \bjet \eta distribution,
    \protect\subref{fig:SR_Whad_Final_l_Whad_particle:bjet1phi} leading \bjet \phi distribution,
    \protect\subref{fig:SR_Whad_Final_l_Whad_particle:bjet2pt} subleading \bjet \pt distribution,
    \protect\subref{fig:SR_Whad_Final_l_Whad_particle:bjet2eta} subleading \bjet \eta distribution,
    \protect\subref{fig:SR_Whad_Final_l_Whad_particle:bjet2phi} subleading \bjet \phi distribution,
    \protect\subref{fig:SR_Whad_Final_l_Whad_particle:mu} pile-up distribution,
    \protect\subref{fig:SR_Whad_Final_l_Whad_particle:njets} number of jets distribution, and
    \protect\subref{fig:SR_Whad_Final_l_Whad_particle:nbjets} number of \bjets distribution.
    The four predictions differ only in the modelling of the \tW
    contributions, where the predictions using dynamic scales are marked
    DR (nominal sample) and DS while the legacy samples with fixed scales
    are marked DRF and DSF.\@
    The ratio panel displays the ratio of the total MC to the data.
    Expected data uncertainty is displayed using error bars.
    Shaded areas show MC statistical uncertainty only.
  }%
  \label{fig:SR_Whad_Final_l_Whad_particle:additional2}
\end{figure}


